% =============================================================================
\section{BioBERT NLP Layer}
% =============================================================================

\begin{frame}{Step 4: Fine-tuned BioBERT overview}
  English text $X$ passes through our fine-tuned BioBERT model $\to$ acuity prediction $\hat{\mathbf{y}}$.

  \vspace{0.3em}
  \begin{equation*}
    X \xrightarrow{\tau} (t_i) \xrightarrow{E} H^{(0)} \xrightarrow{L \times} H^{(L)} \xrightarrow{\mathrm{head}} \hat{\mathbf{y}}
  \end{equation*}

  \footnotesize
  \begin{tabular}{@{}ll@{}}
    \toprule
    \textbf{Predicted acuity} & \textbf{Maps to SATS} \\
    \midrule
    Level 1--2 & Red / Orange (highly urgent) \\
    Level 3 & Yellow (moderate) \\
    Level 4--5 & Green (non-urgent / routine) \\
    \bottomrule
  \end{tabular}

  \plain{Fine-tuning taught BioBERT our hospital's phrasing. It answers ``how urgent does this sound?'' --- that answer goes to fusion with vitals and Bayesian when needed.}
\end{frame}

\begin{frame}{Step 4: Softmax classification}
  \begin{equation*}
    \hat{\mathbf{y}} = \mathrm{softmax}(W h_{\mathrm{[CLS]}} + \mathbf{b}), \quad
    \hat{y}_c = P(\text{acuity} = c \mid X)
  \end{equation*}

  \examplebox{
    \textbf{Chief complaint:} thunderclap headache, worsening with movement\\
    \begin{tabular}{@{}lrr@{}}
      \toprule
      Level & Meaning & Probability \\
      \midrule
      1 & Most urgent & 1\% \\
      2 & Highly urgent & \textbf{94\%} \\
      3 & Moderate & 3\% \\
      4--5 & Routine & 2\% \\
      \bottomrule
    \end{tabular}\\
    Prediction = Level 2 $\to$ maps to \sats{triageOrange}{Orange}
  }

  \plain{Softmax converts raw scores into percentages over 5 acuity levels --- highest wins.}
\end{frame}

\begin{frame}{Step 4: Confidence gate}
  \begin{equation*}
    \text{confidence} = \max_c \hat{y}_c, \quad
    \text{if confidence} < \tau \Rightarrow \text{flag for Bayesian fallback}
  \end{equation*}
  where $\tau \approx 0.85$.

  \examplebox{
    \textbf{High confidence (passes):} $\max(\hat{\mathbf{y}}) = 0.94 > \tau = 0.85$ $\to$ trust BioBERT.\\[0.3em]
    \textbf{Borderline (flagged):} Patient TG-7OKLDLKAQ --- glaucoma complaint:\\
    Level 2: 85.4\%, Level 3: 14.1\% $\Rightarrow$ confidence $0.854$ --- fusion may call Bayesian fallback.
  }

  \plain{When the model is unsure, we do not silently trust it --- fusion escalates to Bayesian reasoning.}
\end{frame}

\begin{frame}{BioBERT: inside the model --- embeddings}
  \begin{equation*}
    E(t_i) = E_{\mathrm{word}}(t_i) + E_{\mathrm{pos}}(i) + E_{\mathrm{seg}}(t_i)
  \end{equation*}

  Each token gets three vectors summed: word identity, position in sentence, segment (question vs.\ answer).

  \examplebox{
    \textbf{Input:} ``crushing chest pain radiating to left arm''\\
    Tokens: [crush, \#\#ing, chest, pain, radiate, \#\#ing, left, arm]\\
    Each token $\to$ 768-dim vector after embedding layer.
  }
\end{frame}

\begin{frame}{BioBERT: self-attention}
  \begin{equation*}
    \mathrm{Attention}(Q, K, V) = \mathrm{softmax}\!\left(\frac{QK^\top}{\sqrt{d_k}}\right) V
  \end{equation*}

  \examplebox{
    \textbf{Query token ``crushing''} attends most strongly to:
    \begin{itemize}
      \item ``crushing'' --- 0.38
      \item ``chest'' --- 0.35
      \item ``pain'' --- 0.22
    \end{itemize}
    Clinical co-occurrence drives urgency detection.
  }
\end{frame}

\begin{frame}{BioBERT: transformer layers and [CLS]}
  \begin{align*}
    Z^{(\ell)} &= \mathrm{LayerNorm}\!\Bigl(H^{(\ell-1)} + \mathrm{MultiHead}(H^{(\ell-1)})\Bigr) \\
    H^{(\ell)} &= \mathrm{LayerNorm}\!\Bigl(Z^{(\ell)} + \mathrm{FFN}(Z^{(\ell)})\Bigr)
  \end{align*}

  \vspace{0.3em}
  After $L=12$ layers:
  \begin{equation*}
    h_{\mathrm{[CLS]}} = H^{(L)}_{1,:} \in \mathbb{R}^{768}
    \qquad
    \hat{\mathbf{y}} = \mathrm{softmax}(W h_{\mathrm{[CLS]}} + \mathbf{b})
  \end{equation*}

  \examplebox{
    ``crushing chest pain'' $\to$ Level 2: 91\%, Level 3: 6\%, others: 3\%\\
    Confidence $0.91 \to$ \sats{triageOrange}{Orange} at fusion.
  }
\end{frame}

\begin{frame}{BioBERT: how it was trained --- pre-training}
  \begin{equation*}
    \mathcal{L}_{\mathrm{MLM}} = -\sum_{i \in \mathcal{M}} \log P(t_i \mid H^{(L)}; \theta)
  \end{equation*}

  \textbf{Masked Language Modelling} on PubMed medical English teaches BioBERT clinical vocabulary.

  \examplebox{
    ``The patient presented with severe \underline{\hspace{1cm}} pain''\\
    Masked token predicted: \textbf{chest}\\
    Model learns co-occurrence of medical terms before ever seeing triage labels.
  }
\end{frame}

\begin{frame}{BioBERT: fine-tuning on hospital data}
  \begin{equation*}
    \mathcal{L} = -\sum_{i=1}^{N} \log P(y_i \mid X_i; \theta)
  \end{equation*}

  Minimise cross-entropy vs.\ nurse-assigned acuity labels on $\sim$80,000 chief complaints.
  \textbf{Vitals excluded} --- model learns urgency from language alone.

  \examplebox{
    Row: ``contraception advice, intermittent'' $\to$ nurse label Level 5\\
    $P(\text{Level 5} \mid X) = 0.97 \Rightarrow \mathcal{L} \mathrel{+}= -\log(0.97) \approx 0.03$
  }

  \footnotesize
  \begin{tabular}{@{}ll@{}}
    \toprule
    \textbf{Setting} & \textbf{Value} \\
    \midrule
    Base model & \texttt{Yuvrajxms09/biobert-triage-classifier} \\
    Data merge & train.csv $\bowtie$ chief\_complaints.csv ($\sim$80k rows) \\
    Epochs & 3 (13,500 steps) \\
    Best eval\_loss & \textbf{0.001812} \\
    \bottomrule
  \end{tabular}
\end{frame}

\begin{frame}{Data exploration and findings}
  Training and inference on \textbf{$\sim$100,000} patient records from fine-tuned BioBERT CSVs.

  \footnotesize
  \begin{columns}[T]
    \column{0.48\textwidth}
    \textbf{Dataset sizes}
    \begin{tabular}{@{}ll@{}}
      \toprule
      File & Rows \\
      \midrule
      train.csv & 80,000 \\
      test.csv & 20,000 \\
      chief\_complaints.csv & 100,000 \\
      triage\_predictions\_results.csv & 100,000 \\
      \bottomrule
    \end{tabular}

    \column{0.48\textwidth}
    \textbf{Acuity distribution (train)}
    \begin{tabular}{@{}lrr@{}}
      \toprule
      Level & Count & \% \\
      \midrule
      1 & 3,222 & 4.0\% \\
      2 & 13,439 & 16.8\% \\
      3 & 28,921 & 36.2\% \\
      4 & 23,020 & 28.8\% \\
      5 & 11,398 & 14.2\% \\
      \bottomrule
    \end{tabular}
  \end{columns}

  \vspace{0.4em}
  \textbf{Confidence:} 99.6\% of predictions at confidence 1.0; \textbf{361} rows at 0.6--0.99 (Bayesian candidates).

  \plain{Some urgent cases are overconfidently misclassified --- fusion with TEWS/Bayesian mitigates this.}
\end{frame}

\begin{frame}{Example predictions from CSV}
  \footnotesize
  \begin{tabular}{@{}p{2.2cm}p{4.5cm}ll@{}}
    \toprule
    \textbf{Patient ID} & \textbf{Complaint} & \textbf{Pred.} & \textbf{Conf.} \\
    \midrule
    TG-UXRGA9UCO & thunderclap headache, worsening with movement & Level 2 & 1.00 \\
    TG-B19DBBS2G & contraception advice, intermittent & Level 5 & 1.00 \\
    TG-7OKLDLKAQ & acute angle closure glaucoma in known patient & Level 2 & 0.854 \\
    TG-2MIQ3DLS7 & acute angle closure glaucoma with nausea & Level 2 & 0.855 \\
    \bottomrule
  \end{tabular}

  \vspace{0.5em}
  \normalsize
  Glaucoma cases at $\sim$0.85 confidence illustrate why the $\tau \approx 0.85$ gate and Bayesian fallback exist.
\end{frame}
